% GENERATED FILE. Edit canonical lessons, then run npm run generate:books.
% canonical-source-hash: fnv1a64:d0dea4a6934a23f4
% canonical-lessons: ML-C410-sucippikkuka, ML-C410-unnipparayuka, ML-C410-cuntikkanikkuka, ML-C410-prasaktamaya, ML-C410-aprasaktamaya

\chapter{Making a Point}
\label{ch:ml-making-a-point}
\hlchaptermodality{410}

\begin{chapteropening}
\textbf{By the end of this chapter:} \emph{I can say to indicate, to emphasise, to point out, relevant and irrelevant.}

Complete the last lesson of chapter 410: I can say to indicate, to emphasise, to point out, relevant and irrelevant.
\end{chapteropening}

\section[\texorpdfstring{\ml{സൂചിപ്പിക്കുക} (sūcippikkuka)}{സൂചിപ്പിക്കുക (sūcippikkuka)}]{\ml{സൂചിപ്പിക്കുക} (sūcippikkuka) — to indicate; to hint}

\label{lesson:ML-C410-sucippikkuka}



\begin{quote}
Before the new one: say the Malayalam for impartial, then the Malayalam for an explanation.
\end{quote}

\subsection*{You'll want to know: \ml{സൂചിപ്പിക്കുക}}
\textbf{\ml{സൂചിപ്പിക്കുക}} — \emph{sūcippikkuka} — \textquotedblleft{}to indicate; to hint\textquotedblright{}.

\begin{quote}
\ml{നാളെ} \ml{മഴ} \ml{വരുമെന്ന്} \ml{റിപ്പോർട്ട്} \ml{സൂചിപ്പിച്ചു}.
\end{quote}

\begin{quote}
\emph{nāḷe maḻa varumennŭ ṟippōrṭṭŭ sūcippiccu.}
\end{quote}

\begin{quote}
The report indicated that rain would come tomorrow.
\end{quote}

Past: \ml{സൂചിപ്പിച്ചു}, \emph{sūcippiccu}.

The noun is \ml{സൂചന}, a hint, a sign.

\begin{cousinweb}
\textbf{\ml{സൂചിപ്പിക്കുക}} (\emph{sūcippikkuka}) is built from Sanskrit \emph{sūcana}, \textquotedblleft{}pointing out, indicating\textquotedblright{} (Monier-Williams), with the Malayalam verb ending, here -\ml{ിപ്പിക്കുക}.
\end{cousinweb}

\begin{grammarlens}[title={What you've built}]
One more word for giving reasons and opinions.
\end{grammarlens}

\subsection*{Your turn}
\begin{itemize}
\raggedright
  \item \emph{Say it:} \emph{sūcippikkuka}
  \item \emph{Say it:} \emph{sūcippikkuka}, once more
  \item \emph{Say it:} \emph{nāḷe maḻa varumennŭ ṟippōrṭṭŭ sūcippiccu.} — The report indicated that rain would come tomorrow.
  \item \emph{Recall:} say the Malayalam for impartial, then the Malayalam for an explanation, then say \emph{sūcippikkuka} again
\end{itemize}

\begin{tcolorbox}[breakable,colback=teal!4,colframe=teal!35!black,title={Before you move on}]
What does \ml{സൂചിപ്പിക്കുക} mean? (\textquotedblleft{}To indicate; to hint\textquotedblright{}.) Say it once more.
\end{tcolorbox}
\section[\texorpdfstring{\ml{ഊന്നിപ്പറയുക} (ūnnippaṟayuka)}{ഊന്നിപ്പറയുക (ūnnippaṟayuka)}]{\ml{ഊന്നിപ്പറയുക} (ūnnippaṟayuka) — to emphasise; to stress}

\label{lesson:ML-C410-unnipparayuka}



\begin{quote}
Before the new one: say the Malayalam for an explanation, then the Malayalam for to indicate.
\end{quote}

\subsection*{You'll want to know: \ml{ഊന്നിപ്പറയുക}}
\textbf{\ml{ഊന്നിപ്പറയുക}} — \emph{ūnnippaṟayuka} — \textquotedblleft{}to emphasise; to stress\textquotedblright{}.

\begin{quote}
\ml{സമയം} \ml{പ്രധാനമാണെന്ന്} \ml{അധ്യാപകൻ} \ml{ഊന്നിപ്പറഞ്ഞു}.
\end{quote}

\begin{quote}
\emph{samayaṁ pradhānamāṇennŭ adhyāpakan ūnnippaṟaññu.}
\end{quote}

\begin{quote}
The teacher stressed that time is important.
\end{quote}

Past: \ml{ഊന്നിപ്പറഞ്ഞു}, \emph{ūnnippaṟaññu}.

Its past is \ml{ഊന്നിപ്പറഞ്ഞു}, the past of \ml{പറയുക} with \ml{ഊന്നി} in front.

\begin{cousinweb}
\textbf{\ml{ഊന്നിപ്പറയുക}} (\emph{ūnnippaṟayuka}) is a Malayalam compound of native Dravidian verbs: \ml{ഊന്നുക}, \textquotedblleft{}to be fixed; to lean, to rest on\textquotedblright{} (DEDR 763), and \ml{പറയുക}, \textquotedblleft{}to say\textquotedblright{} (DEDR 4031). To emphasise is to lean on a word as you say it.
\end{cousinweb}

\begin{grammarlens}[title={What you've built}]
One more word for giving reasons and opinions.
\end{grammarlens}

\subsection*{Your turn}
\begin{itemize}
\raggedright
  \item \emph{Say it:} \emph{ūnnippaṟayuka}
  \item \emph{Say it:} \emph{ūnnippaṟayuka}, once more
  \item \emph{Say it:} \emph{samayaṁ pradhānamāṇennŭ adhyāpakan ūnnippaṟaññu.} — The teacher stressed that time is important.
  \item \emph{Recall:} say the Malayalam for an explanation, then the Malayalam for to indicate, then say \emph{ūnnippaṟayuka} again
\end{itemize}

\begin{tcolorbox}[breakable,colback=teal!4,colframe=teal!35!black,title={Before you move on}]
What does \ml{ഊന്നിപ്പറയുക} mean? (\textquotedblleft{}To emphasise; to stress\textquotedblright{}.) Say it once more.
\end{tcolorbox}
\section[\texorpdfstring{\ml{ചൂണ്ടിക്കാണിക്കുക} (cūṇṭikkāṇikkuka)}{ചൂണ്ടിക്കാണിക്കുക (cūṇṭikkāṇikkuka)}]{\ml{ചൂണ്ടിക്കാണിക്കുക} (cūṇṭikkāṇikkuka) — to highlight; to point out}

\label{lesson:ML-C410-cuntikkanikkuka}



\begin{quote}
Before the new one: say the Malayalam for to indicate, then the Malayalam for to emphasise.
\end{quote}

\subsection*{You'll want to know: \ml{ചൂണ്ടിക്കാണിക്കുക}}
\textbf{\ml{ചൂണ്ടിക്കാണിക്കുക}} — \emph{cūṇṭikkāṇikkuka} — \textquotedblleft{}to highlight; to point out\textquotedblright{}.

\begin{quote}
\ml{റിപ്പോർട്ടിലെ} \ml{ഒരു} \ml{തെറ്റ്} \ml{അവൾ} \ml{ചൂണ്ടിക്കാണിച്ചു}.
\end{quote}

\begin{quote}
\emph{ṟippōrṭṭile oru teṟṟŭ avaḷ cūṇṭikkāṇiccu.}
\end{quote}

\begin{quote}
She pointed out a mistake in the report.
\end{quote}

Past: \ml{ചൂണ്ടിക്കാണിച്ചു}, \emph{cūṇṭikkāṇiccu}.

Also \ml{ചൂണ്ടിക്കാട്ടുക}, with the same meaning.

\begin{cousinweb}
\textbf{\ml{ചൂണ്ടിക്കാണിക്കുക}} (\emph{cūṇṭikkāṇikkuka}) is a Malayalam compound of native Dravidian verbs: \ml{ചൂണ്ടുക}, \textquotedblleft{}to point at, to aim\textquotedblright{} (DEDR 2658), and \ml{കാണിക്കുക}, \textquotedblleft{}to show\textquotedblright{} (DEDR 1443): to point at and show.
\end{cousinweb}

\begin{grammarlens}[title={What you've built}]
One more word for giving reasons and opinions.
\end{grammarlens}

\subsection*{Your turn}
\begin{itemize}
\raggedright
  \item \emph{Say it:} \emph{cūṇṭikkāṇikkuka}
  \item \emph{Say it:} \emph{cūṇṭikkāṇikkuka}, once more
  \item \emph{Say it:} \emph{ṟippōrṭṭile oru teṟṟŭ avaḷ cūṇṭikkāṇiccu.} — She pointed out a mistake in the report.
  \item \emph{Recall:} say the Malayalam for to indicate, then the Malayalam for to emphasise, then say \emph{cūṇṭikkāṇikkuka} again
\end{itemize}

\begin{tcolorbox}[breakable,colback=teal!4,colframe=teal!35!black,title={Before you move on}]
What does \ml{ചൂണ്ടിക്കാണിക്കുക} mean? (\textquotedblleft{}To highlight; to point out\textquotedblright{}.) Say it once more.
\end{tcolorbox}
\section[\texorpdfstring{\ml{പ്രസക്തമായ} (prasaktamāya)}{പ്രസക്തമായ (prasaktamāya)}]{\ml{പ്രസക്തമായ} (prasaktamāya) — relevant}

\label{lesson:ML-C410-prasaktamaya}



\begin{quote}
Before the new one: say the Malayalam for to emphasise, then the Malayalam for to highlight.
\end{quote}

\subsection*{You'll want to know: \ml{പ്രസക്തമായ}}
\textbf{\ml{പ്രസക്തമായ}} — \emph{prasaktamāya} — \textquotedblleft{}relevant\textquotedblright{}.

\begin{quote}
\ml{അത്} \ml{വളരെ} \ml{പ്രസക്തമായ} \ml{ഒരു} \ml{ചോദ്യമാണ്}.
\end{quote}

\begin{quote}
\emph{atŭ vaḷare prasaktamāya oru cōdyamāṇŭ.}
\end{quote}

\begin{quote}
That is a very relevant question.
\end{quote}

The noun is \ml{പ്രസക്തി}, relevance.

\begin{cousinweb}
\textbf{\ml{പ്രസക്തമായ}} (\emph{prasaktamāya}) is built from Sanskrit \emph{prasakta}, \textquotedblleft{}attached to, connected with\textquotedblright{} (Monier-Williams), with the adjective ending -\ml{ആയ}: what is relevant is attached to the matter in hand.
\end{cousinweb}

\begin{grammarlens}[title={What you've built}]
One more word for giving reasons and opinions.
\end{grammarlens}

\subsection*{Your turn}
\begin{itemize}
\raggedright
  \item \emph{Say it:} \emph{prasaktamāya}
  \item \emph{Say it:} \emph{prasaktamāya}, once more
  \item \emph{Say it:} \emph{atŭ vaḷare prasaktamāya oru cōdyamāṇŭ.} — That is a very relevant question.
  \item \emph{Recall:} say the Malayalam for to emphasise, then the Malayalam for to highlight, then say \emph{prasaktamāya} again
\end{itemize}

\begin{tcolorbox}[breakable,colback=teal!4,colframe=teal!35!black,title={Before you move on}]
What does \ml{പ്രസക്തമായ} mean? (\textquotedblleft{}Relevant\textquotedblright{}.) Say it once more.
\end{tcolorbox}
\section[\texorpdfstring{\ml{അപ്രസക്തമായ} (aprasaktamāya)}{അപ്രസക്തമായ (aprasaktamāya)}]{\ml{അപ്രസക്തമായ} (aprasaktamāya) — irrelevant}

\label{lesson:ML-C410-aprasaktamaya}



\begin{quote}
Before the new one: say the Malayalam for to highlight, then the Malayalam for relevant.
\end{quote}

\subsection*{You'll want to know: \ml{അപ്രസക്തമായ}}
\textbf{\ml{അപ്രസക്തമായ}} — \emph{aprasaktamāya} — \textquotedblleft{}irrelevant\textquotedblright{}.

\begin{quote}
\ml{ആ} \ml{ചോദ്യം} \ml{ഇപ്പോൾ} \ml{അപ്രസക്തമാണ്}.
\end{quote}

\begin{quote}
\emph{ā cōdyaṁ ippōḷ aprasaktamāṇŭ.}
\end{quote}

\begin{quote}
That question is irrelevant now.
\end{quote}

The Sanskrit prefix \ml{അ}- makes an opposite: \ml{പ്രസക്തമായ}, \ml{അപ്രസക്തമായ}.

\begin{cousinweb}
\textbf{\ml{അപ്രസക്തമായ}} (\emph{aprasaktamāya}) is built from Sanskrit \emph{a-}, \textquotedblleft{}not\textquotedblright{}, and \emph{prasakta}, \textquotedblleft{}attached to, connected with\textquotedblright{} (Monier-Williams), with the adjective ending -\ml{ആയ}.
\end{cousinweb}

\begin{grammarlens}[title={What you've built}]
One more word for giving reasons and opinions.
\end{grammarlens}

\subsection*{Your turn}
\begin{itemize}
\raggedright
  \item \emph{Say it:} \emph{aprasaktamāya}
  \item \emph{Say it:} \emph{aprasaktamāya}, once more
  \item \emph{Say it:} \emph{ā cōdyaṁ ippōḷ aprasaktamāṇŭ.} — That question is irrelevant now.
  \item \emph{Recall:} say the Malayalam for to highlight, then the Malayalam for relevant, then say \emph{aprasaktamāya} again
\end{itemize}

\begin{tcolorbox}[breakable,colback=teal!4,colframe=teal!35!black,title={Before you move on}]
What does \ml{അപ്രസക്തമായ} mean? (\textquotedblleft{}Irrelevant\textquotedblright{}.) Say it once more.
\end{tcolorbox}
